\begin{tikzpicture}[
    every node/.style={font=\sffamily},
    st/.style={rounded corners=4pt, draw=violet!50, fill=violet!8,
               minimum height=0.9cm, minimum width=2.3cm, text centered,
               font=\sffamily\small\ttfamily},
    fin/.style={st, draw=violet!65!black, line width=1pt, fill=violet!14},
    tr/.style={-{Stealth[length=6pt]}, thick, gray!55!black},
    tl/.style={font=\sffamily\footnotesize, text=gray!45!black, fill=white,
               inner sep=1.5pt, align=center},
]
% ── Estado inicial ──
\fill[violet!65!black] (-6.7,0) circle (2.5pt);
% ── Estados ──
\node[st] (pending)   at (-4.8, 0)   {pending};
\node[st] (running)   at (-0.4, 0)   {running};
\node[st] (resetting) at (-0.4, 2.6) {resetting};
\node[st] (stopping)  at (-0.4, -2.6){stopping};
\node[fin] (completed) at (5.0, 2.0) {completed};
\node[fin] (timeout)   at (5.0, 0.5) {timeout};
\node[fin] (failed)    at (5.0, -1.0){failed};
\node[fin] (stopped)   at (5.0, -2.6){stopped};

% ── Transiciones ──
\draw[tr] (-6.7,0) -- (pending);
\draw[tr] (pending) -- (running) node[tl, pos=0.5] {\texttt{run start} /\\provisionado};
% running → terminales (derecha)
\draw[tr] (running) -- (completed) node[tl, pos=0.62] {targets\\cumplidos};
\draw[tr] (running) -- (timeout)   node[tl, pos=0.7]  {límite de\\tiempo};
\draw[tr] (running) -- (failed)    node[tl, pos=0.62] {error};
% running → stopping → stopped
\draw[tr] (running) -- (stopping) node[tl, pos=0.5] {stop del\\educador};
\draw[tr] (stopping) -- (stopped);
% running ⇄ resetting (reset: dos flechas paralelas)
\draw[tr] ([xshift=-9pt]running.north) -- ([xshift=-9pt]resetting.south)
    node[tl, pos=0.5, left=1pt] {\texttt{run reset}};
\draw[tr] ([xshift=9pt]resetting.south) -- ([xshift=9pt]running.north)
    node[tl, pos=0.5, right=1pt] {nuevo\\\texttt{env\_id}};
\end{tikzpicture}
